\uniterablesection{Висновки}

У результаті виконання курсової роботи спроєктовано й розроблено інформаційну систему моніторингу особистих звичок користувачів, яка поєднує ручне ведення записів з автоматичним збором даних із зовнішніх сервісів. Поставлену мету досягнуто, а всі визначені на початку завдання — розв'язано.

У процесі роботи отримано такі результати:

\begin{enumerate}[noitemsep, topsep=2pt, leftmargin=*, label=\arabic*)]
    \item проаналізовано предметну область трекінгу звичок, розглянуто механізми формування звичок та досліджено наявні рішення; виявлено, що жодне з них не поєднує гнучкість довільних звичок з автоматичним збором активності розробника;
    \item на основі аналізу сформульовано функціональні та нефункціональні вимоги до системи;
    \item обґрунтовано вибір архітектури модульного моноліту та технологічного стеку на платформах .NET 10 і Angular;
    \item спроєктовано доменну модель за принципами предметно-орієнтованого проєктування, структуру бази даних PostgreSQL та REST-інтерфейс системи;
    \item реалізовано серверну частину з поділом на обмежені контексти, розділенням команд і запитів (CQRS) та подієвою взаємодією між модулями;
    \item реалізовано механізм автоматичного періодичного опитування сервісів GitHub і Strava, що самостійно створює записи звичок на основі виявленої активності;
    \item розроблено клієнтський вебзастосунок з наочною візуалізацією прогресу у вигляді теплової карти та стрічки активності;
    \item проведено модульне й функціональне тестування, яке підтвердило працездатність ключових сценаріїв, зокрема автоматичного формування записів за реальними подіями GitHub.
\end{enumerate}

Головна практична цінність роботи полягає в усуненні основної вади традиційних трекерів звичок — залежності від дисциплінованості користувача в момент обліку. Завдяки автоматичному збору активності частина записів формується без участі людини, що знижує навантаження на користувача й підвищує достовірність даних.

Закладені архітектурні рішення залишають широкий простір для розвитку системи. Подальші напрями вдосконалення охоплюють: додавання нових інтеграцій (наприклад, із сервісами навчання чи відстеження сну), розширення аналітики та формування рекомендацій, реалізацію сповіщень і нагадувань, а також розроблення мобільного клієнта. Завдяки слабкому зв'язуванню модулів і подієвій взаємодії ці розширення можна впроваджувати без істотного переписування наявного коду.
